\chapter*{Abstract}
\addcontentsline{toc}{chapter}{Abstract}

Deep learning has transformed artificial intelligence, yet traditional architectures are constrained to data with regular Euclidean structure (sequences, grids, voxels). Geometric Deep Learning (GDL) emerges as a unifying paradigm that extends these techniques to non-Euclidean domains---graphs, manifolds, groups---through the systematic incorporation of geometric symmetries. This book presents a rigorous mathematical analysis of GDL organized around the ``5~Gs'' framework: Geometry, Groups, Graphs, Geodesics, and Gauges.

The development begins with the formalization of classical neural networks using Lebesgue measure theory, including a complete proof of Cybenko's universal approximation theorem. The GDL framework is built around the five Gs: \emph{Geometry} provides the structure of data domains through Riemannian manifolds, tangent spaces, and curvature; \emph{Groups} formalize symmetries via Lie groups and algebras, representation theory (Peter-Weyl and Schur theorems), and the equivariance principle; \emph{Graphs} discretize continuous domains through spectral analysis and the graph Laplacian; \emph{Geodesics} provide Riemannian metrics, the exponential map, and parallel transport for information propagation; and \emph{Gauges} introduce local invariances through principal bundles and connections. This mathematical apparatus is applied to the formalization of equivariant convolutional networks---over finite groups, compact and non-compact Lie groups, and manifolds---and geometric attention mechanisms, including equivariant Transformers and hyperbolic attention.

Through formal proofs and detailed analyses, this book shows how fundamental geometric principles---equivariance, invariance, locality---unify convolutional networks, graph networks, and Transformers under a common geometric message passing framework. The result is a view of GDL as a mathematical framework that not only explains existing architectures but guides the design of new ones.

\vspace{0.5cm}
\begin{flushleft}
	\textbf{Keywords}: Geometric Deep Learning, Group Theory, Equivariance, Riemannian Manifolds, Representation Theory, Lie Groups, Convolutional Networks, Graph Neural Networks, Transformers, Message Passing
\end{flushleft}
